\section{Introduction}
\label{sec:intro}

Gait analysis plays a crucial role in the diagnosis and management of neurodevelopmental disorders in children, particularly Cerebral Palsy (CP)~\cite{gage2009treatment}. The Edinburgh Visual Gait Score (EVGS)~\cite{read2003edinburgh} is a widely adopted clinical tool that enables systematic assessment of gait deviations through observational scoring of 17 biomechanical items per limb. However, EVGS scoring requires trained clinicians and is inherently subjective, motivating the development of automated computational approaches. Similarly, classifying gait patterns in children with Bilateral Spastic CP---such as True Equinus, Jump Gait, Apparent Equinus, and Crouch Gait~\cite{winters1987gait, rodda2004classification}---remains challenging due to subtle inter-class differences and the rarity of certain patterns.

The CVPR 2026 AI for Children Challenge (CV4CHL)~\cite{cv4chl2026} introduces the Children Gait Pose Sequence (CGPS) dataset, comprising 2D keypoint annotations across 339,236 frames from 110 pediatric patients captured in multi-view clinical settings. The challenge features two tracks: Track~1 requires predicting 17 binary EVGS items per limb, while Track~2 involves classifying five gait patterns for patients with Bilateral Spastic CP. The final ranking is determined by the average score across both tracks.

This challenge presents several unique difficulties that distinguish it from standard skeleton-based action recognition. First, the dataset is extremely small by machine learning standards---78 training patients for Track~1 and only 18 for Track~2---making overfitting a dominant concern. Second, the test data exhibits substantially higher keypoint noise than the training data, creating a significant distribution shift that renders conventional cross-validation unreliable. Third, certain EVGS items and gait patterns are severely imbalanced (\eg, knee hyperextension appears in only 7.4\% of left limbs; Crouch Gait has only 2 training samples). Fourth, the clinical nature of the task demands that predictions align with biomechanical plausibility, not merely statistical patterns.

In this report, we present a domain-knowledge-guided hybrid framework that addresses these challenges through five key contributions:

\begin{itemize}[leftmargin=*,nosep]
    \item A \textbf{hierarchical data cleaning pipeline} incorporating confidence-based keypoint filtering, double-pass temporal smoothing, and per-joint outlier detection, which we identify as the single most impactful factor for leaderboard performance. 
    \item A \textbf{hybrid clinical-ML prediction strategy} for Track~2 that fuses ensemble machine learning predictions with rule-based clinical heuristics at the decision layer, achieving a 40\% rule weighting that outperforms both pure ML and pure rule-based approaches.
    \item \textbf{Cross-track knowledge transfer} that uses Track~1 EVGS predictions as discriminative features for Track~2 gait classification, exploiting the strong correlation between EVGS profiles and gait patterns.
    \item \textbf{Patient-level classification with bilateral symmetry priors} for Track~2, leveraging the clinical observation that 91\% of training patients exhibit identical gait patterns in both limbs.
    \item A \textbf{cross-version submission ensemble} that combines predictions from models with diverse feature engineering pipelines, providing robustness through methodological diversity rather than model diversity alone.
\end{itemize}

Through systematic ablation across 14 solution iterations, we derive 12 empirical laws governing the behavior of machine learning on small clinical datasets. These findings reveal counter-intuitive phenomena---such as cross-validation being anti-correlated with test performance, and naive combination of individually strong methods degrading overall accuracy---that we believe offer valuable insights for the broader medical AI community.
